\section*{Bloque A. Fundamentos teóricos y conceptuales}

\noindent\textit{Preguntas 1 a 30. Evalúan la comprensión rigurosa del modelo clásico de regresión lineal, estimador MCO, supuestos de Gauss-Markov, Teorema BLUE, inferencia estadística ($t, F$), $R^2$, modelos no lineales, multicolinealidad, heterocedasticidad, autocorrelación y endogeneidad.}

\q{En la representación matricial del Modelo Lineal General $Y = X\beta + u$, el estimador de Mínimos Cuadrados Ordinarios (MCO) $\hat{\beta}$ se obtiene minimizando la suma de residuos al cuadrado y su expresión cerrada es:}
\begin{opts}
    \item $\hat{\beta} = (X'X) X'Y$
    \item $\hat{\beta} = (X'X)^{-1} X'Y$
    \item $\hat{\beta} = (X'Y)^{-1} X'X$
    \item $\hat{\beta} = X'(XX')^{-1} Y$
\end{opts}

\q{El supuesto de 'Exogeneidad Estricta' en el modelo de regresión lineal establece formalmente que $E[u|X] = 0$. La consecuencia teórica directa más importante de este supuesto es que:}
\begin{opts}
    \item el estimador MCO es asintóticamente eficiente pero presenta sesgo en muestras finitas.
    \item los errores estándar de los estimadores convergen a cero cuando el tamaño de muestra es finito.
    \item el coeficiente de determinación $R^2$ alcanza siempre un valor unitario en muestras finitas.
    \item el estimador MCO es insesgado respecto al vector de parámetros poblacionales ($E[\hat{\beta}|X] = \beta$).
\end{opts}

\q{Para que la matriz $(X'X)^{-1}$ exista y el estimador MCO pueda calcularse unívocamente, se requiere el supuesto de:}
\begin{opts}
    \item normalidad estricta de la variable dependiente $Y$ en las observaciones.
    \item homocedasticidad estricta y ausencia de autocorrelación serial en los errores.
    \item rango de columna completo de la matriz $X$, descartando multicolinealidad perfecta.
    \item independencia temporal completa entre las perturbaciones estocásticas.
\end{opts}

\q{El supuesto de perturbaciones esféricas en el modelo clásico de Gauss-Markov abarca dos condiciones fundamentales sobre la matriz de varianzas-covarianzas condicional del error $Var(u|X) = \sigma^2 I_n$:}
\begin{opts}
    \item estacionariedad en sentido estricto y ergodicidad asintótica del término de error.
    \item normalidad bivariada, asimetría nula y curtosis exacta igual a tres en los residuos.
    \item homocedasticidad (varianza constante $\sigma^2$) y ausencia de autocorrelación serial.
    \item independencia estocástica lineal entre los regresores y la variable dependiente medida.
\end{opts}

\q{El célebre Teorema de Gauss-Markov demuestra que, bajo los supuestos clásicos de regresión lineal, el estimador MCO:}
\begin{opts}
    \item coincide con el estimador de máxima verosimilitud siempre que los errores sigan distribución uniforme.
    \item alcanza un coeficiente $R^2$ superior al de cualquier estimador no lineal bajo errores esféricos.
    \item es invariante en magnitud frente a cualquier transformación no lineal en la escala de las variables.
    \item es el Mejor Estimador Lineal Insesgado (MELI), con varianza mínima entre estimadores lineales insesgados.
\end{opts}

\q{El estimador insesgado de la varianza poblacional de las perturbaciones estocásticas ($\sigma^2$) en una regresión múltiple con $n$ observaciones y $k$ parámetros (incluyendo la constante) es:}
\begin{opts}
    \item $s^2 = \frac{\sum_{i=1}^n e_i^2}{n - 1}$
    \item $s^2 = \frac{\sum_{i=1}^n e_i}{n - k}$
    \item $s^2 = \frac{e'e}{n - k}$
    \item $s^2 = \frac{\sum_{i=1}^n e_i^2}{n}$
\end{opts}

\q{El coeficiente de determinación $R^2$ mide la proporción de la variación total de la variable dependiente explicada por el modelo de regresión. Una propiedad matemática intrínseca del $R^2$ es que:}
\begin{opts}
    \item disminuye necesariamente al introducir cualquier regresor adicional irrelevante.
    \item su valor siempre oscila entre $-1$ y $+1$ cuando el modelo incluye constante.
    \item nunca disminuye al agregar variables explicativas adicionales, aunque no sean significativas.
    \item es igual a cero solo cuando todos los coeficientes estimados son negativos.
\end{opts}

\q{Para penalizar la inclusión innecesaria de regresores que no aportan poder explicativo genuino, se utiliza el coeficiente $R^2$ ajustado ($\bar{R}^2$), el cual se define como:}
\begin{opts}
    \item $\bar{R}^2 = 1 - \frac{SCR}{STC} \sqrt{n - k}$
    \item $\bar{R}^2 = 1 - \frac{n - k}{n - 1} (1 - R^2)$
    \item $\bar{R}^2 = 1 - \frac{SEC / (n - k)}{STC / (n - 1)}$
    \item $\bar{R}^2 = 1 - \frac{SCR / (n - k)}{STC / (n - 1)}$
\end{opts}

\q{Si además de los supuestos de Gauss-Markov se añade el supuesto de Normalidad de los errores ($u|X \sim N(0, \sigma^2 I_n)$), la distribución muestral exacta del estimador MCO es:}
\begin{opts}
    \item $\hat{\beta}|X \sim N(\beta, \sigma^2 (X'X)^{-1})$
    \item $\hat{\beta}|X \sim t_{n-k}(\beta, \sigma^2 I_n)$
    \item $\hat{\beta}|X \sim F_{k, n-k}(0, \sigma^2 I_n)$
    \item $\hat{\beta}|X \sim \chi^2_{n-k}(\beta, \sigma^2 X'X)$
\end{opts}

\q{Para contrastar la hipótesis nula de significancia individual $H_0: \beta_j = 0$ contra $H_1: \beta_j \neq 0$, se utiliza el estadístico $t$ de Student definido como:}
\begin{opts}
    \item $t = \frac{SCR}{SEC}$, que bajo $H_0$ sigue una distribución normal estándar en muestras finitas.
    \item $t = \frac{\hat{\beta}_j - 1}{EE(\hat{\beta}_j)}$, con distribución normal estándar para cualquier tamaño de muestra.
    \item $t = \frac{\hat{\beta}_j}{EE(\hat{\beta}_j)}$, que bajo $H_0$ sigue una $t$ con $n - k$ grados de libertad.
    \item $t = \frac{\hat{\beta}_j}{s^2}$, que bajo $H_0$ sigue una distribución $\chi^2$ con $n - k$ grados de libertad.
\end{opts}

\q{La prueba $F$ de significancia conjunta del modelo contrasta la hipótesis nula $H_0: \beta_2 = \beta_3 = \dots = \beta_k = 0$ (todos los coeficientes de las pendientes son cero simultáneamente) frente a la alternativa de que al menos uno es distinto de cero. Su fórmula a partir del $R^2$ es:}
\begin{opts}
    \item $F = t^2 \times \frac{n - k}{k - 1}$, con distribución $t_{n-k}$ bajo $H_0$.
    \item $F = \frac{1 - R^2}{R^2} (n - k)$, con distribución $F_{n-k, k-1}$ bajo $H_0$.
    \item $F = \frac{R^2 / (k - 1)}{(1 - R^2) / (n - k)}$, con distribución $F_{k-1, n-k}$ bajo $H_0$.
    \item $F = \frac{R^2 / (n - k)}{(1 - R^2) / (k - 1)}$, con distribución $\chi^2_{k-1}$ bajo $H_0$.
\end{opts}

\q{En un modelo econométrico con especificación doble logarítmica (log-log) $\ln(Y_i) = \beta_0 + \beta_1 \ln(X_i) + u_i$, el coeficiente de la pendiente $\beta_1$ se interpreta directamente como:}
\begin{opts}
    \item la tasa de crecimiento porcentual instantánea de la constante ante variaciones unitarias en $X$.
    \item la elasticidad constante de $Y$ respecto a $X$ (variación porcentual de $Y$ por $1\%$ en $X$).
    \item la variación absoluta en niveles de $Y$ ante una variación unitaria en el logaritmo de $X$.
    \item la semielasticidad de $Y$ ante un cambio discreto de una unidad física en el regresor $X$.
\end{opts}

\q{En un modelo semilogarítmico del tipo log-lineal $\ln(Y_i) = \beta_0 + \beta_1 X_i + u_i$ (como la ecuación de ingresos de Mincer), el parámetro $\beta_1$ se interpreta como:}
\begin{opts}
    \item una semielasticidad: un aumento unitario en $X$ cambia $Y$ en aproximadamente $100 \times \beta_1\%$.
    \item el cambio absoluto promedio en $Y$ dividido por el logaritmo neperiano de la media de $X$.
    \item el coeficiente de correlación lineal simple ajustado entre las variables originales $Y$ y $X$.
    \item una elasticidad constante unitaria: ante un cambio de $1\%$ en $X$, $Y$ cambia en $\beta_1$ unidades.
\end{opts}

\q{La presencia de 'Multicolinealidad Imperfecta' o severa en una regresión múltiple tiene como consecuencia econométrica que:}
\begin{opts}
    \item los estimadores MCO pierden insesgadez y consistencia en muestras grandes.
    \item MCO sigue siendo insesgado, pero las varianzas de los estimadores aumentan, reduciendo los estadísticos $t$.
    \item la matriz $X'X$ deja de ser invertible para calcular el vector $\hat{\beta}$.
    \item el coeficiente de determinación $R^2$ cae automáticamente a valores negativos.
\end{opts}

\q{El problema de 'Heterocedasticidad' se define formalmente cuando la varianza condicional de los errores no es constante entre observaciones:}
\begin{opts}
    \item $X'u \neq 0$ en muestras pequeñas con constante incluida.
    \item $E[u_i|X] \neq 0$, variando según la media de cada regresor.
    \item $Cov(u_i, u_j|X) \neq 0$ para todo par de observaciones $i \neq j$.
    \item $Var(u_i|X) = \sigma_i^2$, variando entre observaciones.
\end{opts}

\q{Cuando existe heterocedasticidad en los errores pero se aplica MCO convencional:}
\begin{opts}
    \item la prueba de Durbin-Watson rechaza sistemáticamente la hipótesis nula de perturbaciones esféricas.
    \item el estimador de la varianza residual $s^2$ converge a cero a medida que aumenta el número de variables explicativas.
    \item $\hat{\beta}$ sigue siendo insesgado, pero pierde eficiencia y los errores estándar están sesgados, invalidando los tests $t$ y $F$.
    \item los estimadores $\hat{\beta}$ pierden la insesgadez y la consistencia asintótica incluso bajo exogeneidad estricta.
\end{opts}

\q{La prueba general de White para detectar heterocedasticidad consiste en regresar los residuos al cuadrado $e_i^2$ sobre:}
\begin{opts}
    \item el logaritmo de los errores estandarizados y los valores ajustados al cubo.
    \item todos los regresores originales, sus cuadrados y sus productos cruzados no redundantes ($n R^2 \sim \chi^2_p$).
    \item la variable dependiente rezagada un periodo junto con las variables ficticias.
    \item las diferencias de primer orden de las variables explicativas sin constante.
\end{opts}

\q{La prueba de Breusch-Pagan para heterocedasticidad asume que la varianza del error es una función lineal de un conjunto de variables explicativas $Z_1, \dots, Z_m$. Su hipótesis nula $H_0$ es:}
\begin{opts}
    \item $H_0: \sigma_1^2 = \dots = \sigma_n^2$ (los coeficientes de los regresores $Z$ en la auxiliar de $e_i^2$ son nulos).
    \item $H_0$: los residuos siguen una distribución exponencial negativa con varianza estrictamente decreciente.
    \item $H_0$: el término constante de la regresión original es cero y la media de los residuos supera la varianza.
    \item $H_0$: no existe multicolinealidad severa ni autocorrelación de primer orden entre los regresores $X$.
\end{opts}

\q{El método de Mínimos Cuadrados Ponderados (WLS / MCG) para corregir la heterocedasticidad cuando se conoce la estructura de la varianza $Var(u_i|X) = \sigma^2 h_i$ consiste en:}
\begin{opts}
    \item eliminar el $10\%$ de las observaciones con mayores residuos al cuadrado para evitar apalancamiento extremo.
    \item elevar al cuadrado todas las variables explicativas y estimar la regresión auxiliar sin término constante.
    \item dividir cada observación de todas las variables (incluida la constante) entre $\sqrt{h_i}$ y aplicar MCO.
    \item aplicar primeras diferencias logarítmicas a la variable dependiente para estabilizar la escala residual.
\end{opts}

\q{Los errores estándar robustos a la heterocedasticidad propuestos por White (estimador 'Sandwich' o HC0):}
\begin{opts}
    \item exigen que el tamaño de muestra sea menor a 30 observaciones y corrigen la media de la variable dependiente.
    \item mantienen intactos los $\hat{\beta}$ de MCO pero corrigen la matriz de covarianzas mediante el estimador tipo sándwich.
    \item requieren que las perturbaciones sigan obligatoriamente una distribución $t$ de Student con varianza conocida.
    \item modifican las estimaciones puntuales de los coeficientes $\hat{\beta}$ mediante mínimos cuadrados ponderados en dos etapas.
\end{opts}

\q{El fenómeno de 'Autocorrelación Serial' de primer orden $AR(1)$ en modelos de series de tiempo implica que la perturbación en el periodo $t$ sigue el proceso:}
\begin{opts}
    \item $u_t = \sum_{i=1}^t X_i$, sin término de innovación aleatorio ni memoria.
    \item $u_t = \rho u_{t-1} + \varepsilon_t$, con $|\rho| < 1$ y $\varepsilon_t$ ruido blanco.
    \item $u_t = u_t^2 + \varepsilon_t$, donde la varianza depende linealmente del intercepto.
    \item $u_t = \alpha Y_t + \varepsilon_t$, con $\alpha > 1$ y $\varepsilon_t$ tendencia determinística.
\end{opts}

\q{Cuando existe autocorrelación serial positiva ($\rho > 0$) en los errores de un modelo de series de tiempo estimado por MCO:}
\begin{opts}
    \item el estadístico de Durbin-Watson converge sistemáticamente hacia un valor cercano a 4 en muestras grandes.
    \item los regresores se vuelven perfectamente colineales con sus rezagos, haciendo singular la matriz $(X'X)$.
    \item los estimadores de los coeficientes $\hat{\beta}$ se vuelven sesgados e inconsistentes aun con regresores estrictamente exógenos.
    \item $\hat{\beta}$ sigue insesgado bajo exogeneidad estricta, pero los errores estándar subestiman la varianza, inflando los tests $t$.
\end{opts}

\q{El estadístico $d$ de Durbin-Watson para contrastar autocorrelación serial de primer orden se aproxima algebraicamente a la fórmula $d \approx 2(1 - \hat{\rho})$. Por ende, en ausencia total de autocorrelación ($\hat{\rho} = 0$), el valor del estadístico es:}
\begin{opts}
    \item $d \approx 2$
    \item $d \approx 0$
    \item $d \approx -1$
    \item $d \approx 4$
\end{opts}

\q{Una limitación crucial del estadístico $d$ de Durbin-Watson es que queda estrictamente invalidado si:}
\begin{opts}
    \item los datos provienen de encuestas de corte transversal ordenadas de forma estrictamente aleatoria.
    \item el modelo contiene un término constante junto con más de dos regresores continuos exógenos.
    \item el número de observaciones de la serie supera el umbral de 100 periodos trimestrales.
    \item el modelo incluye la variable dependiente rezagada ($Y_{t-1}$) entre los regresores explicativos.
\end{opts}

\q{La prueba de Breusch-Godfrey (o prueba del Multiplicador de Lagrange - LM) para autocorrelación es superior a la de Durbin-Watson porque:}
\begin{opts}
    \item contrasta autocorrelación de orden superior $AR(p)$ y es válida aun cuando existan rezagos de $Y_t$.
    \item no requiere calcular los residuos de MCO ni estimar una regresión auxiliar sobre los regresores.
    \item se aplica únicamente a modelos con coeficientes estandarizados sin intercepto ni tendencias.
    \item elimina automáticamente los valores atípicos mediante ponderación iterativa de Cochrane-Orcutt.
\end{opts}

\q{Para realizar inferencia estadística válida en series de tiempo en presencia conjunta de heterocedasticidad y autocorrelación de forma desconocida, se utilizan:}
\begin{opts}
    \item los errores estándar HAC consistentes de Newey y West (1987).
    \item los contrastes de raíz unitaria aumentados de Dickey-Fuller (ADF).
    \item el filtro cíclico de Hodrick-Prescott aplicado sobre los regresores.
    \item los errores estándar clásicos de MCO ajustados por grados de libertad.
\end{opts}

\q{El problema de 'Endogeneidad' se produce cuando uno o más regresores están correlacionados con la perturbación estocástica ($Cov(X, u) \neq 0$). Sus tres causas más comunes en economía son:}
\begin{opts}
    \item transformación logarítmica de la variable dependiente $Y$.
    \item homocedasticidad condicional, normalidad y estacionariedad.
    \item sesgo de variable omitida, error de medición en $X$ y simultaneidad.
    \item ausencia de constante, colinealidad exacta y exceso de muestra.
\end{opts}

\q{La consecuencia econométrica más grave de la endogeneidad en el estimador de MCO es que:}
\begin{opts}
    \item el estimador $\hat{\beta}$ es sesgado e inconsistente ($plim \, \hat{\beta} \neq \beta$).
    \item el coeficiente de correlación de Pearson entre los regresores se vuelve exactamente igual a 1.
    \item el estimador MCO conserva su consistencia asintótica pero pierde eficiencia en muestras grandes.
    \item los errores estándar calculados convergen a cero provocando un sobreajuste de los residuos.
\end{opts}

\q{Para resolver la endogeneidad mediante el método de Variables Instrumentales (VI / IV), una variable instrumental válida $Z$ debe satisfacer dos condiciones indispensables:}
\begin{opts}
    \item estar altamente correlacionada con la perturbación $u$ y no correlacionada con $X$.
    \item relevancia del instrumento ($Cov(Z, X) \neq 0$) y exogeneidad ($Cov(Z, u) = 0$).
    \item seguir una distribución uniforme continua y ser ortogonal a la variable explicativa $X$.
    \item ser una transformación cuadrática exacta de la variable dependiente observada en el modelo.
\end{opts}

\q{El estimador de Mínimos Cuadrados en Dos Etapas (MC2E / 2SLS) resuelve la endogeneidad siguiendo el procedimiento:}
\begin{opts}
    \item Etapa 1: Diferenciar las series temporales; Etapa 2: Aplicar mínimos cuadrados ponderados sobre los residuos.
    \item Etapa 1: Regresar $Y$ sobre los instrumentos $Z$; Etapa 2: Regresar la variable endógena $X$ sobre los residuos.
    \item Etapa 1: Regresar $X$ sobre los instrumentos $Z$ para obtener $\hat{X}$; Etapa 2: Regresar $Y$ sobre $\hat{X}$ por MCO.
    \item Etapa 1: Calcular la matriz de correlaciones parciales; Etapa 2: Eliminar las variables con menor estadístico $t$.
\end{opts}

